% Original attached proof of night:P1.9; displaced only by exact cited specialization.
\begin{proof}
\begin{enumerate}[label=\arabic*.]
\item Let $\Sigma$ consist of the cones over all simplices of the periodic
Kuhn triangulation at height one, together with their faces.  For an
ordered maximal simplex $v_0,\ldots,v_3$, the columns $(v_j,1)$ form a
unimodular four-by-four matrix.  If $u_j$ are its dual integral basis,
the corresponding chart has coordinates $x_j=\chi^{u_j}$ and is
$\mathbb C^4$.  The base character is the sum of the four $u_j$, so
\[
q=x_0x_1x_2x_3.
\]
On a common face we invert the coordinates of omitted rays and change
dual bases; these are mutually inverse integral monomial maps.  These
explicit charts and overlaps construct a smooth toric space over the
disk with a reduced normal-crossings zero fibre.

\item The period group acts by
\[
(z,q)\longmapsto
\bigl(q^\lambda\exp(2\pi iC(q)\lambda)z,q\bigr),
\qquad \lambda\in\mathbb Z^3.
\]
The monomial part takes the cone over $\sigma$ to the cone over
$\sigma+\lambda$; the unit part is torus multiplication holomorphic in
$q$, which extends over each chart.  These maps compose according to
addition in $\mathbb Z^3$.  On the zero fibre a fixed point would
require a nonzero translation to stabilize a bounded simplex, which
is impossible.  For small nonzero $q$, the imaginary part of the
period matrix is $\operatorname{Im}(s)I_3+\operatorname{Im}C(q)$ and
is invertible; hence the action is free there too.

\item We give a compactness check for the infinite-chart construction.  Put
$u=\log|z|/\log|q|$.  A period changes $u$ by
$(I+E(q))\lambda$, where $E(q)=O(1/|\log|q||)$ uniformly.
On a small disk $(I+E)^{-1}$ has norm at most two.  Reduction of
$(I+E)^{-1}u$ modulo $\mathbb Z^3$ therefore gives a representative
with $u$ in a fixed bounded region.  Only finitely many simplices
meet that region.  In a simplex containing $u$, the numbers
$\log|x_j|/\log|q|$ are its barycentric coordinates; thus
$|x_j|\leq1$.  Finitely many closed polydisks in these charts give
compact representatives above a smaller closed base disk.  Every
central point is included in such representatives by translating its
finite star into the same bounded region and then using its complete
toric component.

\item The action is proper.  For compact coordinate neighborhoods approaching
the central fibre, only finitely many cones and their faces occur;
their translates can meet only when the translation is bounded.
Equivalently, the valuation estimate just given and the uniform bound
on $(I+E)^{-1}$ bound the translations for a sequence with $q\to0$.
On a closed annulus away from zero, properness is the usual
properness of a full lattice action.  These two checks give finitely
many intersecting translates of any pair of compact neighborhoods.
The quotient is therefore Hausdorff, is smooth by freeness, and is
proper over the disk by the compact representatives.  The map is flat
because $q$ is a nonzero divisor in each smooth chart.  Exponentiation
$z_j=\exp(2\pi i\zeta_j)$ identifies the nonzero fibres with
$\mathbb C^3/(\mathbb Z^3+Z(q)\mathbb Z^3)$.

\item A $p$-simplex supplies a stratum $(\mathbb C^*)^{3-p}$.
Enumeration of all faces of the six Kuhn tetrahedra, retaining their
translation identifications, gives $(1,7,12,6)$.
Additivity of compactly supported Euler characteristic and
$\chi_c(\mathbb C^*)=0$ leave six point contributions.
Inclusion of an index-$d$ period sublattice induces a finite covering
of these toric quotients.  Each simplex orbit has exactly $d$ lifts,
which proves the stated index dependence by an actual comparison map.

\end{enumerate}
\end{proof}

% Original attached proof of maps:src:F.1.1; displaced only by exact cited specialization.
\begin{proof}
\begin{enumerate}[label=\arabic*.]
\item A fundamental cube has six unimodular tetrahedra.  Taking the
quotient by an index-$d$ sublattice multiplies every cell count
by $d$, so the quotient cell counts are $d(1,7,12,6)$.
A cell of dimension $r$ corresponds to a torus stratum
$(\mathbb C^*)^{3-r}$.  Only the $6d$ point strata contribute
to Euler characteristic.  Euler integration over the regular
base gives zero since $\chi(T^6)=0$, and a finite free torus
quotient has Euler characteristic $0/|G|=0$.  Additivity now
gives the formula.  This checks the rank-three instance of the
localization mechanism in \ExternalPaperRef{proposition~}{4.1}{VI} and \ExternalPaperRef{proposition~}{4.3}{VI}.  It requires the specified
toroidal filling, not just its unipotent monodromy matrix.

\end{enumerate}
\end{proof}
